\begin{center}
Problema G

Comportamento Dinâmico

{\small Nome do arquivo: G.c, G.cpp, G.java, ou G.py}

{\large Timelimit: $1$}

{\tiny Por Heverton Barros de Macêdo}

\end{center}

\footnotesize

	Durante uma investigação sobre o comportamento de sistemas dinâmicos, Wolfram definiu o seguinte cenário para efetuar experimentos: dada uma cadeia de $N$ células, linearmente distribuídas, cada célula armazena um único estado (bit), e além disso as células das extremidades são vizinhas entre si. Para cada instante de tempo $t$ a evolução dessa cadeia ocorre simultaneamente em todas as células, considerando uma regra de evolução previamente definida. Essa regra é aplicada a cada vizinhança de $K=3$ células gerando um novo estado, conforme indicado na Tabela~\ref{tab:evolucao}.
	
	\begin{table}[h!]
		\caption{Regra aplicada para evolução do experimento de Wolfram.}
		\label{tab:evolucao}
		
		\centering
		\footnotesize{
		\begin{tabular}{|c|c|c|c|c|c|c|c|c|}
			\hline
			\cellcolor[HTML]{DDDDDD} Regra para vizinhança & 000 & 001 & \cellcolor[HTML]{FFAAAA} 010 & 011 & 100 & 101 & 110 & 111 \\ \hline
			\cellcolor[HTML]{DDDDDD} Novo estado para célula central & 1 & 0 &  \cellcolor[HTML]{AAFFAA} 0 & 1 & 0 & 0 & 0 & 1 \\ \hline
		\end{tabular}
		}
	\end{table}

Como exemplo, considere um experimento (veja a Tabela~\ref{tab:exemplo_evolucao}) com a cadeia de bits formada por $11$ células ($N = 11$), em que no instante $t = 0$ apenas um bit, exatamente no centro da cadeia, possui o estado $1$ e as demais células possuem o estado $0$, evoluindo por $4$ passos de tempo ($t = 4$).
%Como exemplo, considere um experimento (veja a Tabela~\ref{tab:exemplo_evolucao}) com a cadeia de bits no espaço linear formada por $11$ células ($N = 11$), em que apenas um bit, exatamente no centro da cadeia ($t = 0$), possui o estado $1$ e as demais células possuem o estado $0$, evoluindo por $4$ passos de tempo.

A Tabela~\ref{tab:exemplo_evolucao} ilustra o experimento com a configuração inicial ($t = 0$) formada por: $00000100000$. No próximo passo de tempo ($t = 1$), todas as células evoluem (podem alterar o estado) considerando a vizinhança indicada na regra (Tabela~\ref{tab:evolucao}). Em destaque na cor vermelha (Tabelas~\ref{tab:evolucao} e ~\ref{tab:exemplo_evolucao}) está a atualização da célula central, formada pela vizinhança $010$ que, segundo a regra, deve evoluir no próximo passo de tempo ($t = 1$) para o estado $0$ (destaque na cor verde). É possível notar no exemplo, a evolução da cadeia inicial por $4$ passos de tempo. Vale destacar que as células na extremidade (esquerda e direta) estão conectadas (formato de anel - toroidal)

	\begin{table}[h!]
	\caption{Exemplo de um experimento por $4$ passos de tempo.}
	\label{tab:exemplo_evolucao}
	
	\centering
	\footnotesize{	
	\begin{tabular}{|c|c|c|c|c|c|c|c|c|c|c|c|}
		\hline
		Tempo & \multicolumn{11}{c|}{Cadeia linear de N células}\\ \hline
    	$t = 0$ & 0 & 0 & 0 & 0 & \cellcolor[HTML]{FFAAAA} 0 & \cellcolor[HTML]{FFAAAA} 1 & \cellcolor[HTML]{FFAAAA} 0 & 0 & 0 & 0 & 0 \\ \hline
		$t = 1$ & 1 & 1 & 1 & 1 &  0 & \cellcolor[HTML]{AAFFAA} 0 & 0 & 1 & 1 & 1 & 1 \\ \hline
		$t = 2$ & 1 & 1 & 1 & 0 &  1 & 0 & 0 & 1 & 1 & 1 & 1 \\ \hline
		$t = 3$ & 1 & 1 & 0 & 0 &  0 & 0 & 0 & 1 & 1 & 1 & 1 \\ \hline
		$t = 4$ & 1 & 0 & 0 & 1 &  1 & 1 & 0 & 1 & 1 & 1 & 1 \\ \hline
	\end{tabular}
	}	
\end{table}

%\pagebreak
Sua missão é ajudar Wolfram a obter a configuração dos estados após $t$ passos de tempo, considerando uma
sequência inicial qualquer no tempo $t = 0$ e a aplicação da regra indicada no enunciado.

\vspace{0.3cm}
\textbf{Entrada}

A entrada contém vários casos de teste. Cada linha contém um inteiro $t$ ($0~<~t~\leq~1000$), que representa
a quantidade de passos de evolução, seguido pela cadeia de bits no tempo $t = 0$, que pode variar conforme a
quantidade de células da amostra $N$ ($3 \leq N < 100$).

\vspace{0.3cm}
\textbf{Saída}

O programa deve produzir como saída a configuração da cadeia no espaço linear de células após a evolução
de $t$ passos de tempo.

\begin{table}[h]
    \centering
    \footnotesize
    \begin{tabular}{|p{8cm}|p{8cm}|}
        \hline
        \begin{center}
            \vspace{-0.3cm}
            \textbf{Exemplos de Entrada}
            \vspace{-0.3cm}
        \end{center} 
         &  
        \begin{center}
            \vspace{-0.3cm}
            \textbf{Exemplos de Saída}
            \vspace{-0.3cm}
        \end{center} \\ \hline
		    \texttt{1 00000100000} & \texttt{11110001111} \\
         \hline
   		    \texttt{2 00000100000} & \texttt{11100101111} \\
         \hline    
   		    \texttt{3 00000100000} & \texttt{11000001111} \\
         \hline    
   		    \texttt{578 00000100000} & \texttt{00000111101} \\
         \hline    
   		    \texttt{978 011001111000110111010011} & \texttt{100111000001110000011101} \\
         \hline    
    \end{tabular}
    \caption{Exemplos de entradas e saídas}    
\end{table}

\normalsize

\newpage

\begin{center}
\noindent
\lstinputlisting{abobora_25/G/G5.cpp}
\end{center}